\section{工程剖面、真实代码量与演进纪律}

\begin{keyidea}
学习项目的规模用进入目标系统的真实机制衡量，不用测试、工具或文档行数制造
体量。测试和教材仍是必要交付物，但它们分别证明和解释代码，不计入操作系统
本体代码量。
\end{keyidea}

\subsection{一、仓库中的五类产物}

\begin{longtable}{L{0.17\linewidth}L{0.24\linewidth}L{0.31\linewidth}L{0.18\linewidth}}
\toprule
目录 & 运行位置 & 职责 & 计入真实代码 \\
\midrule
\filepath{source/} & 目标 x86-64 机器 & 固件、基础模块、后续内核与用户程序 & 是 \\
\filepath{tests/} & 宿主机/QEMU 调度 & 单元、集成、随机和整机断言 & 否 \\
\filepath{tools/} & 宿主机 & 构建、审计、镜像和 QEMU 生命周期 & 否 \\
\filepath{books/} & 文档构建环境 & 原理教材、图表和代码走读 & 否 \\
\filepath{web/} & 网站运行环境 & 在线阅读与发布门户 & 否 \\
\bottomrule
\end{longtable}

CMakeLists 和链接脚本属于构建与布局描述。链接脚本对 ROM 地址至关重要，但
处理器不会把它作为指令执行，因此本书把它列为工程资产，不计入目标代码行。
统计口径保留这种边界，避免随着测试和文档扩充产生“内核已经很大”的错觉。

\subsection{二、当前真实代码量}

本书构建时调用 \code{python3} \code{tools/os.py} \code{source-metrics} 自动扫描
\filepath{source/}。只接受 \code{.asm}、\code{.cpp}、\code{.hpp}、
\code{.inc} 和 \code{.tpp}；空行与纯注释行不计入代码行，行尾注释所在的
指令或语句仍计入。

\begin{longtable}{L{0.32\linewidth}L{0.22\linewidth}L{0.36\linewidth}}
\toprule
指标 & 当前值 & 含义 \\
\midrule
真实源文件 & \OsProductionFileCount & 进入目标系统语义的源文件 \\
物理行 & \OsProductionPhysicalLineCount & 包含空行和注释的原始行数 \\
有效代码行 & \OsProductionCodeLineCount & 非空、非纯注释行 \\
freestanding C++20 & \OsProductionCppCodeLineCount & 头文件和源文件有效行 \\
NASM Intel 汇编 & \OsProductionAssemblyCodeLineCount & 固件、模式切换与目标加载器有效行 \\
\bottomrule
\end{longtable}

当前阶段汇编占比较高，因为复位固件、实模式、模式切换、ATA PIO 和第一版
ELF64 加载器都必须在 C++ ABI 建立前工作。v0.4 已经建立 64 位栈和第一次
C++ 交接；从 v0.5 起，高层机制会主要进入 C++20，汇编只保留入口、异常桩、
上下文切换和必要硬件边界。这个变化反映系统层次演进，不是追求某种语言比例。

\subsection{三、代码行不是完成度}

\OsProductionCodeLineCount 行有效目标代码已经完成真实复位、串口、ATA PIO、
64 MiB 页表、长模式、两遍 ELF64 装载、BootInfo 和 C++ 内核入口，但尚未
建立内核自己的 IDT、TSS 与异常分发。相反，短小的复位向量只有三个有效字节，
却承担整个系统是否获得执行权的关键边界。完成度必须由机制和验收项表示。

\begin{longtable}{L{0.20\linewidth}L{0.30\linewidth}L{0.40\linewidth}}
\toprule
表面指标 & 能说明的事实 & 不能说明的事实 \\
\midrule
代码行增加 & 仓库出现更多目标语句 & 机制正确、边界完整或可维护 \\
测试行增加 & 覆盖意图和输入数量增加 & 测试断言与规格一致 \\
串口有输出 & CPU 至少执行到某个标记 & 后续状态、错误路径和硬件全部正确 \\
阶段验收通过 & 当前明确契约有重复证据 & 尚未实现的下一阶段能力存在 \\
\bottomrule
\end{longtable}

\subsection{四、自动统计的可重复性}

统计实现位于宿主工具，不进入目标镜像。书籍 Makefile 在检查和 PDF 构建前
重新生成 LaTeX 宏；源代码增加后，表格数字随构建更新。工具本身由单元测试
覆盖 C++ 行注释、块注释、NASM 注释和扩展名过滤。

统计值必须与提交绑定。发布记录保存当时阶段、提交和口径；未来改变统计规则时
先更新规则说明，再重新生成数字，不能把不同口径的曲线直接比较。

\subsection{五、从 v0.2 到 v1.0 的规模增长来源}

后续真实代码增长主要来自异常入口、页分配、更细粒度页表、驱动、调度、同步、
文件系统和用户程序。每个模块都先建立类型与失败语义，再添加正常路径。测试和
教材会同步增长，但继续从真实代码量中排除。

\subsection{六、阅读工程的稳定顺序}

阅读一个阶段先看发布记录和 ADR，确定范围与决策；再看链接脚本和公开头文件，
建立地址与接口；随后走读实现和最终字节；最后看测试如何证伪错误输入。这个
顺序让代码保持在硬件和契约上下文中，不会把某个函数误解为完整机制。
